\documentclass[11pt]{article}
\usepackage[T1]{fontenc}
\usepackage[utf8]{inputenc}
\usepackage{lmodern}
\usepackage[a4paper,margin=28mm]{geometry}
\usepackage{amsmath,amssymb,amsthm,booktabs}
\usepackage{microtype}
\usepackage{url}
\usepackage[hidelinks]{hyperref}
\hypersetup{pdftitle={Nonconvergence of the normalized size of the second remainder set},
  pdfauthor={Xiaojian Zeng}}
\newtheorem{theorem}{Theorem}
\newtheorem{lemma}[theorem]{Lemma}
\newtheorem{corollary}[theorem]{Corollary}
\newcommand{\PP}{\mathcal P}
\DeclareMathOperator{\lcm}{lcm}
\title{Nonconvergence of the normalized size\protect\\of the second remainder set}
\author{Xiaojian Zeng\thanks{PBC School of Finance, Tsinghua University,
100190, China. Email: \href{mailto:xiaojianzeng@foxmail.com}{xiaojianzeng@foxmail.com}.}}
\date{}
\begin{document}
\maketitle

\begin{abstract}
Starting with $S_0(n)=\{1,\ldots,\lfloor n/2\rfloor\}$, define
$S_{j+1}(n)=\{n\bmod k:k\in S_j(n),\ k>0\}$ and $s_j(n)=|S_j(n)|$.
We prove that $s_2(n)/n$ does not converge, settling the second-iterate
case of Problem 2 of Baraskar and Vukusic. An exact congruence
description of positive second remainders gives a finite-cutoff count
on each arithmetic progression. Two such counts, evaluated by exhaustive
residue enumeration and prime-power divisibility, give rationally
separated subsequence bounds with a positive gap exceeding $0.0098$.
Explicit error terms also give separated bounds for every progression
parameter $t\ge203$. No result for higher iterates or exact global
limit inferior and limit superior is asserted.
\end{abstract}

\section{Statement and context}

For every positive integer $n$, let
\begin{equation}\label{eq:sets}
 \begin{gathered}
 S_0(n)=\{1,\ldots,\lfloor n/2\rfloor\},\qquad s_j(n)=|S_j(n)|,\\
 S_{j+1}(n)=\{n\bmod k:k\in S_j(n),\ k>0\}\quad(j\ge0).
 \end{gathered}
\end{equation}
Repeated remainders are counted once. Zero is retained as a value and
excluded as a modulus. In particular, all these sets are empty for $n=1$.

Baraskar and Vukusic studied the iterated sets in
\eqref{eq:sets}. Their Theorem 1 proves convergence of $s_1(n)/n$;
their Theorem 3 gives, for the second iterate
\cite[Theorems 1 and 3]{BaraskarVukusic2026},
\[
 \frac1{24}\le\liminf_{n\to\infty}\frac{s_2(n)}n
 \le\limsup_{n\to\infty}\frac{s_2(n)}n\le\frac14
 .
\]
Problem 2 in their Section 7 asks for nonconvergence of $s_j(n)/n$
for each fixed $j\ge2$ \cite[p.~16]{BaraskarVukusic2026}.
We prove precisely its $j=2$ case. The proof supplies two explicit
arithmetic progressions whose limiting bounds are strictly separated;
the general bounds above do not give such a separation.

Write $\PP$ for the set of primes and, for an integer $Q\ge3$, define
\begin{equation}\label{eq:period}
 L_Q=\lcm\{pq:3\le q\le Q,\ p\in\PP,\ p<q\}.
\end{equation}
The constants used below are
\begin{align*}
 L&=L_{10}=75600, &
 M&=L_{50}=691374732709966591241323368000,\\
 \alpha&=\frac{30093341768093808664597}{276977102566985955634500}, &
 \beta&=\frac{1182869}{9979200}.
\end{align*}
Their difference is
\begin{equation}\label{eq:delta}
 \delta:=\beta-\alpha
 =\frac{131410063319607383019139}{13294900923215325870456000}>0.
\end{equation}

\begin{theorem}\label{thm:main}
For every integer $t\ge1$,
\begin{align}
 \frac{s_2(Mt)}{Mt}&\le\alpha+\frac{M+3}{Mt},\label{eq:upper}\\
 \frac{s_2(Lt+1)}{Lt+1}&\ge\beta-\frac{L+2}{Lt+1}.\label{eq:lower}
\end{align}
Consequently,
\begin{equation}\label{eq:separation}
 \limsup_{t\to\infty}\frac{s_2(Mt)}{Mt}
 \le\alpha<\beta\le
 \liminf_{t\to\infty}\frac{s_2(Lt+1)}{Lt+1},
\end{equation}
so $s_2(n)/n$ has no limit as $n\to\infty$.
\end{theorem}

The argument uses Euclidean division, finite residue counting, and the
Chinese remainder theorem. The latter identifies residues modulo a
product of relatively prime prime powers with the corresponding tuples
of local residues \cite[Theorem 2.6]{Shoup2009}. We use it only for
exact finite proportions; no probabilistic model for integer
factorization is assumed.

\section{Congruences and a uniform cutoff}

\begin{lemma}\label{lem:membership}
For positive integers $n$ and $r$, one has $r\in S_2(n)$ if and only if
there are an integer $q\ge3$ and a prime $p<q$ such that
\begin{equation}\label{eq:membership}
 (q+1)r\le n-q,\qquad pq\mid(q-1)n+r.
\end{equation}
\end{lemma}
\begin{proof}
Suppose $r\in S_2(n)$ and $r>0$. Choose $k\in S_1(n)$ with
$n=qk+r$ and $0<r<k$. There are integers $a\ge2$ and
$1\le h\le\lfloor n/2\rfloor$ with $n=ah+k$ and $0<k<h$.
Thus $n\ge2(k+1)+k$, which implies $q\ge3$. Choose a prime divisor
$p$ of $a$. Since $(n-k)/p\ge h>k$,
\[
 p<\frac nk-1=q-1+\frac rk<q.
\]
Also $q(n-k)=(q-1)n+r$, giving the divisibility in
\eqref{eq:membership}, and $k\ge r+1$ gives its inequality.

Conversely, \eqref{eq:membership} implies $q\mid n-r$. Set
$k=(n-r)/q$. Then $k\ge r+1$, and cancelling $q$ in
$pq\mid q(n-k)$ gives $p\mid n-k$. This cancellation is valid even
when $p\mid q$. Put $h=(n-k)/p$. Since $p\le q-1$ and $r>0$,
\[
 n-k=(q-1)k+r>pk,
\]
so $h>k$. Moreover, $p\ge2$ and $k>0$ give
$h\le\lfloor n/2\rfloor$. The divisions $n=ph+k$ and $n=qk+r$
therefore prove $k\in S_1(n)$ and $r\in S_2(n)$.
\end{proof}

Let $T_Q(n)$ count the positive remainders admitted by
\eqref{eq:membership} with $3\le q\le Q$. A positive remainder
missing from this count has a witness $q\ge Q+1$, hence
\[
 r\le\frac{n-q}{q+1}<\frac{n}{Q+2}.
\]
Accounting separately for the possible zero remainder gives
\begin{equation}\label{eq:tail}
 0\le s_2(n)-T_Q(n)\le\left\lfloor\frac{n}{Q+2}\right\rfloor+1.
\end{equation}

Fix a residue $b$ modulo $L_Q$. For a residue $a$ modulo $L_Q$, let
$w_b(a)$ be the least $q\in\{3,\ldots,Q\}$ for which
\begin{equation}\label{eq:first}
 pq\mid(q-1)b+a\quad\text{for some }p\in\PP,\ p<q,
\end{equation}
and set $w_b(a)=\infty$ if no such $q$ exists. Define
\begin{equation}\label{eq:coefficient}
 A_q(b)=\#\{a\bmod L_Q:w_b(a)=q\},\qquad
 C_Q(b)=\sum_{q=3}^{Q}\frac{A_q(b)}{L_Q(q+1)}.
\end{equation}

\begin{lemma}\label{lem:count}
If $n\ge Q$ and $n\equiv b\pmod{L_Q}$, then
\begin{equation}\label{eq:count}
 |T_Q(n)-nC_Q(b)|\le L_Q+2.
\end{equation}
In particular,
\begin{equation}\label{eq:bounds}
 C_Q(b)-\frac{L_Q+2}{n}\le\frac{s_2(n)}n
 \le C_Q(b)+\frac1{Q+2}+\frac{L_Q+3}{n}.
\end{equation}
\end{lemma}
\begin{proof}
Every modulus $pq$ in \eqref{eq:first} divides $L_Q$, so the
congruence tests depend only on $a$ and $b$. Their endpoints
\[
 H_q(n)=\left\lfloor\frac{n-q}{q+1}\right\rfloor
       =\left\lfloor\frac{n+1}{q+1}\right\rfloor-1
\]
are nonincreasing in $q$. Thus a class with $w_b(a)=q$ contributes
exactly its positive representatives up to $H_q(n)$.

Represent the class by $c\in\{1,\ldots,L_Q\}$, using $c=L_Q$ for
zero. Since $n\ge Q$, the endpoint is nonnegative. For an integer
$H\ge0$, the number of these representatives up to $H$ is
\[
 \left\lfloor\frac{H+L_Q-c}{L_Q}\right\rfloor,
\]
which differs from $H/L_Q$ by at most one. Also
$|H_q(n)-n/(q+1)|<2$. The error for one eligible class is therefore
at most $1+2/L_Q$. Summing over at most $L_Q$ classes proves
\eqref{eq:count}. Combining this with \eqref{eq:tail} proves
\eqref{eq:bounds}.
\end{proof}

\section{The two finite coefficients}

\subsection{The residue class one}

For $Q=10$ and $b=1$, testing \eqref{eq:first} on the full set
$\{0,\ldots,75599\}$ gives
\begin{center}
\begin{tabular}{c*{8}{r}}
\toprule
$q$ &3&4&5&6&7&8&9&10\\
\midrule
$A_q(1)$ &12600&12600&6930&7770&4110&3540&165&1128\\
\bottomrule
\end{tabular}
\end{center}
For clarity, the following complete integer enumeration produces this
table. Quotients are processed in increasing order, and a residue is
assigned only at its first occurrence.
\begin{verbatim}
L = 75600
first = [0] * L
counts = [0] * 11
for q in range(3, 11):
    for p in (2, 3, 5, 7):
        if p >= q:
            continue
        m = p*q
        for a in range((1-q) % m, L, m):
            if first[a] == 0:
                first[a] = q
                counts[q] += 1
\end{verbatim}
The loop covers each allowed congruence throughout a complete period.
Consequently,
\begin{equation}\label{eq:beta}
 C_{10}(1)=\frac1{75600}\left(
 \frac{12600}4+\frac{12600}5+\frac{6930}6+\frac{7770}7+
 \frac{4110}8+\frac{3540}9+\frac{165}{10}+\frac{1128}{11}
 \right)=\beta.
\end{equation}

\subsection{The residue class zero}

For $b=0$, the condition is simply $pq\mid r$. Its least quotient
has the following description.

\begin{lemma}\label{lem:firstzero}
For a positive integer $r$, the least integer $q\ge3$ for which
$pq\mid r$ for some prime $p<q$, when it exists, is the minimum of:
the second smallest distinct prime divisor of $r$, if present; and
the numbers $p^2$ with $p\in\PP$ and $p^3\mid r$.
For $w_0$, candidates exceeding $Q$ are discarded.
\end{lemma}
\begin{proof}
If a witness quotient $q$ is prime, the second smallest distinct prime
divisor of $r$ is at most $q$. If $q$ is composite with a prime divisor
$\ell\ne p$, then $q'=\max(p,\ell)<q$ and $p'=\min(p,\ell)$
give an earlier witness. Otherwise $q=p^a$ for an integer $a\ge2$;
then $p^3\mid r$ and $q'=p^2\le q$ is a witness. Conversely, every
listed candidate directly gives a witness.
\end{proof}

Put
\[
 P_x=\prod_{\substack{a<x\\a\in\PP}}\left(1-\frac1a\right),
 \qquad d_q=\frac{A_q(0)}{L_Q}\quad(3\le q\le Q).
\]
Every prime at most $Q$ divides $L_Q$, and $p^3\mid L_Q$ whenever
$p^2\le Q$, by the pair $(q,p)=(p^2,p)$ in \eqref{eq:period}.
Thus the Chinese remainder theorem gives exact product proportions
for all the prime-power conditions used below. Larger prime powers in
$L_Q$ just repeat these proportions. For the zero residue we may use
the positive representative $L_Q$.

For a prime first quotient $\ell$, exactly one prime $p<\ell$ divides
$r$, and $\ell\mid r$. If $p^2<\ell$, we must also exclude
$p^3\mid r$; if $p^2>\ell$, there is no restriction on its exponent.
All other primes below $\ell$ are absent. The density factors for
the distinguished prime $p$ are respectively $1/p-1/p^3$ and $1/p$.
Dividing out its factor $1-1/p$ in $P_\ell$ gives
\begin{equation}\label{eq:primedensity}
 d_\ell=\frac{P_\ell}{\ell}
 \sum_{\substack{p<\ell\\p\in\PP}}h(p,\ell),\qquad
 h(p,\ell)=
 \begin{cases}
  1/p+1/p^2,&p^2<\ell,\\
  1/(p-1),&p^2>\ell.
 \end{cases}
\end{equation}
The alternatives are disjoint, and equality $p^2=\ell$ is impossible.
For a square first quotient $p^2$, its cube must divide $r$ and every
other prime below $p^2$ must be absent. Any such other prime would
give an earlier prime pair. Hence
\begin{equation}\label{eq:squaredensity}
 d_{p^2}=\frac{1}{p^3}\frac{P_{p^2}}{1-1/p}
        =\frac{P_{p^2}}{p^2(p-1)}.
\end{equation}
All remaining $d_q$ are zero by Lemma~\ref{lem:firstzero}.

At $Q=50$, the nonzero densities from
\eqref{eq:primedensity}--\eqref{eq:squaredensity} are shown in
Table~\ref{tab:densities}. Their exact weighted sum is
\begin{equation}\label{eq:c0}
 C_{50}(0)=\sum_{q=3}^{50}\frac{d_q}{q+1}
 =\frac{6191714756605250456368}{69244275641746488908625},
 \qquad C_{50}(0)+\frac1{52}=\alpha.
\end{equation}
\begin{table}[tb]
\centering
\small
\begin{tabular}{rr@{\qquad}rr}
\toprule
$q$&$d_q$&$q$&$d_q$\\
\midrule
3&$1/6$&23&$31104/2187185$\\
4&$1/12$&25&$27648/16900975$\\
5&$1/12$&29&$59242752/5391411025$\\
7&$2/35$&31&$1689209856/167133741775$\\
9&$4/315$&37&$31905792/3829070245$\\
11&$116/3465$&41&$28931530752/3900644404195$\\
13&$16/585$&43&$160298237952/22952212862579$\\
17&$304/15015$&47&$18484212400128/2928046583754721$\\
19&$12224/692835$&49&$38220595200/81013146586099$\\
\bottomrule
\end{tabular}
\caption{All nonzero first-quotient proportions modulo $L_{50}$ for
the residue class $b=0$. Every entry is an exact rational number.}
\label{tab:densities}
\end{table}

\section{Separation, verification, and scope}

\begin{proof}[Proof of Theorem~\ref{thm:main}]
Apply the upper bound in Lemma~\ref{lem:count} with
$(Q,b,n)=(50,0,Mt)$ and use \eqref{eq:c0}. Apply its lower bound
with $(Q,b,n)=(10,1,Lt+1)$ and use \eqref{eq:beta}. For every
$t\ge1$, both integers exceed their respective cutoffs, so this proves
\eqref{eq:upper}--\eqref{eq:lower}. Letting $t$ tend to infinity
and using \eqref{eq:delta} gives \eqref{eq:separation}.
The two index sequences are strictly increasing and unbounded. If
$s_2(n)/n$ converged, both subsequences would have the same limit,
contradicting \eqref{eq:separation}.
\end{proof}

\begin{corollary}\label{cor:finite}
Set $A=(M+3)/M$ and $B=(L+2)/L$. For every integer $t\ge203$,
\[
 \frac{s_2(Mt)}{Mt}\le\alpha+\frac{A}{203}
 <\beta-\frac{B}{203}\le\frac{s_2(Lt+1)}{Lt+1}.
\]
\end{corollary}
\begin{proof}
Use $(L+2)/(Lt+1)\le B/t$ in \eqref{eq:lower} and $A/t\le A/203$
in \eqref{eq:upper}. The gap between the two fixed bounds is
\begin{equation}\label{eq:finitegap}
 \delta-\frac{A+B}{203}
 =\frac{559612671694976819161613011}
 {17543633842515402252748580463000}>0.
\end{equation}
\end{proof}

The threshold $203$ is an explicit sufficient choice. The same coarse
error allowance fails at $202$; no minimality for the actual counts is
claimed. We do not identify the global limit inferior or limit superior.
The cases $j\ge3$ in the original problem are not resolved here.

All displayed arithmetic can be reproduced with the accompanying
\texttt{certificate.py}, using Python~3.10 or later and only the standard
library. It enumerates all $75600$ residues for $C_{10}(1)$, evaluates
\eqref{eq:primedensity}--\eqref{eq:squaredensity} using rational
arithmetic, and checks the two periods and both positive gaps.
No enumeration of all $L_{50}$ residues is required. Running
\texttt{python3 certificate.py --full} also performs the following
inclusion--exclusion cross-check.

As a separate check of the zero-class densities, the arithmetic
verification starts with the raw clauses $pq\mid r$. If $U$ is the
indicator of the union already processed, adding a clause $E$ replaces
it by $U+1_E-U1_E$. Products of divisibility indicators are represented
by least common multiples. This is the finite inclusion--exclusion
identity \cite[Theorem 8.1]{Shoup2009}. For a period divisible by every
modulus, a term $c\,1_{d\mid r}$ contributes exactly $cL_Q/d$
residues. Differences of successive unions recover the first-quotient
counts, providing a check that does not use the prime-square
classification. A separately written verifier also checked
Lemma~\ref{lem:membership} against literal remainder sets for every
$1\le n\le768$ and checked both coefficient calculations.

The use of exact predicates, complete finite coverage, and a separate
implementation has a methodological parallel in Section 6 of
Zeng et al.\ \cite{ZengEtAl2026Certificates}, where finite geometric
enumerations are checked from their defining integer data. Here the
finite computations supply the periodic coefficients; the proofs of
Lemmas~\ref{lem:membership} and~\ref{lem:count} justify the claims
for arbitrary $n$. Small-instance tests alone would not do so.

\paragraph{Computational and AI assistance.}
OpenAI Codex assisted with proof development, literature retrieval,
exact-arithmetic code, a separate mathematical review pass, and
manuscript preparation. The review pass is not external peer review.
The mathematical argument and explicitly specified integer and rational
calculations, rather than language-model output or numerical fitting,
support the conclusions. No proof assistant was used.

\bibliographystyle{plainurl}
\bibliography{references}
\end{document}
